% Owner: Thuan. 150--250 words (currently ~210). Numbers: exps/scf/results/summary.md, exps/qscf/results/summary.md
\begin{abstract}
Exemplar-free class-incremental learning (CIL) methods built on pre-trained transformers store a mean and a
covariance for every class seen so far, yet use these statistics only to retrain the classifier and discard them
at prediction time. We propose Statistics-Calibrated Fusion (SCF), which keeps the trained classifier and adds a
second score computed from the stored statistics, rescaled by a label-free factor so that a single fusion weight
carries over across scores and datasets. The second score can be classical, a closed-form ridge classifier, or
quantum, a pretty-good measurement (PGM) over class density matrices built from the same statistics. On four
benchmarks with three seeds, SCF raises the final accuracy of the RSIAT baseline by 0.38 to 0.67 points, without
retraining the network or storing any data. Every quantum component is compared with a classical twin that
receives the same inputs: the PGM matches ridge ($-0.04$ points, 95\% CI $[-0.17,+0.08]$, 12 paired runs), as our
invariance results and the least-squares view of the PGM predict. The gain comes from complementary errors: the
classifier and the PGM disagree on 14\% of ImageNet-R test images, and fusion raises accuracy on both old and new classes, whereas
training-time changes trade one for the other. A pre-registered $2\times2$ experiment confirms this: SCF adds 0.38
and 0.44 points of average accuracy on classically and quantum-trained networks, while quantum training adds
nothing measurable once SCF is used.
\keywords{Class-incremental learning \and Exemplar-free \and Score fusion \and
Quantum state discrimination \and Pre-trained models}
\end{abstract}
